\section{Learning Objectives}

After completing this chapter, readers will be able to:
\begin{itemize}
    \item Design robust experimentation infrastructure
    \item Implement randomization and assignment systems
    \item Build reliable logging and data collection pipelines
    \item Integrate A/B testing into applications
    \item Debug common implementation issues
\end{itemize}

\section{Experimentation Infrastructure}

A robust A/B testing system requires careful technical implementation.

\subsection{Core Components}

\begin{enumerate}
    \item \textbf{Randomization Service}: Assigns users to variants
    \item \textbf{Configuration Management}: Stores experiment definitions
    \item \textbf{Logging System}: Records assignments and events
    \item \textbf{Analysis Pipeline}: Computes metrics and statistical tests
    \item \textbf{Monitoring Dashboard}: Tracks experiment health
\end{enumerate}

\subsection{Architecture Patterns}

\textbf{Client-side randomization:}
\begin{itemize}
    \item Randomization happens in browser/app
    \item Lightweight, no server dependency
    \item Risk of manipulation
\end{itemize}

\textbf{Server-side randomization:}
\begin{itemize}
    \item Centralized assignment service
    \item More secure and consistent
    \item Requires infrastructure
\end{itemize}

\textbf{Hybrid approach:}
\begin{itemize}
    \item Server assigns, client implements
    \item Balance of benefits
\end{itemize}

\section{Randomization Implementation}

\subsection{Hash-Based Randomization}

Use deterministic hashing for consistent assignments across sessions.

\lstset{language=Python, caption=Hash-Based User Assignment}
\begin{lstlisting}
import hashlib

def assign_variant(user_id, experiment_id, num_variants=2,
                  traffic_allocation=1.0, salt=''):
    """
    Assign user to variant using deterministic hashing.

    Parameters:
    -----------
    user_id : str
        Unique user identifier
    experiment_id : str
        Unique experiment identifier
    num_variants : int
        Number of variants
    traffic_allocation : float
        Fraction of users to include (0.0-1.0)
    salt : str
        Random salt for security

    Returns:
    --------
    variant : int or None
        Assigned variant (0 to num_variants-1) or None if not allocated
    """
    # Create hash input
    hash_input = f"{user_id}:{experiment_id}:{salt}".encode('utf-8')

    # Generate hash
    hash_output = hashlib.md5(hash_input).hexdigest()

    # Convert to integer and normalize to [0, 1]
    hash_int = int(hash_output, 16)
    hash_fraction = (hash_int % 10000) / 10000.0

    # Check traffic allocation
    if hash_fraction >= traffic_allocation:
        return None  # User not in experiment

    # Assign to variant
    variant_fraction = hash_fraction / traffic_allocation
    variant = int(variant_fraction * num_variants)

    return variant

# Example usage
user_id = "user_12345"
experiment_id = "button_color_test"

# Consistent assignment
for _ in range(3):
    variant = assign_variant(user_id, experiment_id, num_variants=2)
    print(f"User {user_id} assigned to variant: {variant}")

# Different experiments give different assignments
for exp_id in ["exp_1", "exp_2", "exp_3"]:
    variant = assign_variant(user_id, exp_id)
    print(f"Experiment {exp_id}: variant {variant}")
\end{lstlisting}

\textbf{Key properties:}
\begin{itemize}
    \item Deterministic: same user always gets same variant
    \item Independent: different experiments independent
    \item Uniform: variants equally likely
    \item Secure: hard to game with salt
\end{itemize}

\subsection{Handling Edge Cases}

\begin{enumerate}
    \item \textbf{New vs. returning users:} Ensure consistent assignment

    \item \textbf{Logged out users:} Use cookies with careful handling

    \item \textbf{Multi-device users:} Decision: device-level or user-level?

    \item \textbf{Bots and crawlers:} Filter before assignment

    \item \textbf{Internal traffic:} Exclude or separate analysis
\end{enumerate}

\section{Configuration Management}

Store experiment definitions in a centralized system.

\lstset{language=Python, caption=Experiment Configuration Schema}
\begin{lstlisting}
import json
from datetime import datetime

class ExperimentConfig:
    """Experiment configuration manager."""

    def __init__(self, config_path):
        with open(config_path, 'r') as f:
            self.config = json.load(f)

    def get_experiment(self, experiment_id):
        """Retrieve experiment configuration."""
        return self.config.get('experiments', {}).get(experiment_id)

    def is_active(self, experiment_id):
        """Check if experiment is currently active."""
        exp = self.get_experiment(experiment_id)
        if not exp:
            return False

        now = datetime.now()
        start = datetime.fromisoformat(exp.get('start_date'))
        end = datetime.fromisoformat(exp.get('end_date'))

        return start <= now <= end and exp.get('status') == 'active'

# Example configuration file (JSON)
example_config = {
    "experiments": {
        "checkout_button_test": {
            "name": "Checkout Button Color Test",
            "description": "Testing red vs blue checkout button",
            "owner": "product_team@example.com",
            "start_date": "2025-01-01T00:00:00",
            "end_date": "2025-01-31T23:59:59",
            "status": "active",
            "variants": [
                {"id": 0, "name": "control", "allocation": 0.5},
                {"id": 1, "name": "treatment", "allocation": 0.5}
            ],
            "metrics": {
                "primary": "conversion_rate",
                "secondary": ["revenue_per_user", "bounce_rate"],
                "guardrails": ["page_load_time", "error_rate"]
            },
            "targeting": {
                "countries": ["US", "CA", "UK"],
                "platforms": ["web", "mobile_web"]
            }
        }
    }
}
\end{lstlisting}

\section{Logging and Data Collection}

\subsection{Event Schema}

Design a comprehensive event schema capturing all necessary data.

\lstset{language=Python, caption=Event Logging System}
\begin{lstlisting}
import uuid
from datetime import datetime
import json

class ExperimentLogger:
    """Log experiment assignments and events."""

    def __init__(self, storage_backend):
        self.storage = storage_backend

    def log_assignment(self, user_id, experiment_id, variant,
                      timestamp=None, metadata=None):
        """Log user assignment to experiment variant."""
        event = {
            'event_id': str(uuid.uuid4()),
            'event_type': 'assignment',
            'timestamp': timestamp or datetime.utcnow().isoformat(),
            'user_id': user_id,
            'experiment_id': experiment_id,
            'variant': variant,
            'metadata': metadata or {}
        }

        self.storage.write(event)
        return event

    def log_conversion(self, user_id, experiment_id, conversion_type,
                      value=None, timestamp=None, metadata=None):
        """Log conversion event."""
        event = {
            'event_id': str(uuid.uuid4()),
            'event_type': 'conversion',
            'timestamp': timestamp or datetime.utcnow().isoformat(),
            'user_id': user_id,
            'experiment_id': experiment_id,
            'conversion_type': conversion_type,
            'value': value,
            'metadata': metadata or {}
        }

        self.storage.write(event)
        return event

# Example: In-memory storage for demonstration
class InMemoryStorage:
    def __init__(self):
        self.events = []

    def write(self, event):
        self.events.append(event)

    def query(self, event_type=None, experiment_id=None):
        results = self.events
        if event_type:
            results = [e for e in results if e['event_type'] == event_type]
        if experiment_id:
            results = [e for e in results
                      if e.get('experiment_id') == experiment_id]
        return results

# Usage
storage = InMemoryStorage()
logger = ExperimentLogger(storage)

# Log assignment
logger.log_assignment('user_001', 'exp_123', variant=1,
                     metadata={'device': 'mobile'})

# Log conversion
logger.log_conversion('user_001', 'exp_123', 'purchase',
                     value=49.99, metadata={'product_id': 'SKU_789'})

# Query events
assignments = storage.query(event_type='assignment',
                           experiment_id='exp_123')
print(f"Found {len(assignments)} assignment events")
\end{lstlisting}

\subsection{Data Quality Checks}

Implement automated checks:

\begin{enumerate}
    \item \textbf{Sample Ratio Mismatch (SRM):} Verify allocation matches expectations

    \item \textbf{Assignment Consistency:} Users don't switch variants

    \item \textbf{Timestamp Validity:} Events in reasonable order

    \item \textbf{Data Completeness:} No missing critical fields

    \item \textbf{Metric Sanity:} Values within expected ranges
\end{enumerate}

\lstset{language=Python, caption=Data Quality Validation}
\begin{lstlisting}
def validate_experiment_data(assignments_df, events_df):
    """
    Run data quality checks on experiment data.

    Returns:
    --------
    dict : Validation results with issues flagged
    """
    issues = []

    # Check 1: Sample Ratio Mismatch
    variant_counts = assignments_df['variant'].value_counts()
    expected_ratio = 1.0 / len(variant_counts)
    actual_ratios = variant_counts / len(assignments_df)

    from scipy.stats import chisquare
    expected = [len(assignments_df) * expected_ratio] * len(variant_counts)
    chi2, p_value = chisquare(variant_counts.values, expected)

    if p_value < 0.01:
        issues.append({
            'type': 'SRM',
            'severity': 'high',
            'message': f'Sample ratio mismatch detected (p={p_value:.4f})',
            'details': dict(variant_counts)
        })

    # Check 2: Multiple assignments
    multi_assigned = assignments_df.groupby('user_id')['variant'].nunique()
    multi_assigned_users = (multi_assigned > 1).sum()

    if multi_assigned_users > 0:
        issues.append({
            'type': 'consistency',
            'severity': 'high',
            'message': f'{multi_assigned_users} users assigned to multiple variants',
        })

    # Check 3: Events before assignment
    merged = events_df.merge(assignments_df, on='user_id', how='left')
    merged['assignment_timestamp'] = pd.to_datetime(merged['assignment_timestamp'])
    merged['event_timestamp'] = pd.to_datetime(merged['event_timestamp'])

    early_events = (merged['event_timestamp'] <
                   merged['assignment_timestamp']).sum()

    if early_events > 0:
        issues.append({
            'type': 'timing',
            'severity': 'medium',
            'message': f'{early_events} events occurred before assignment'
        })

    return {
        'valid': len(issues) == 0,
        'issues': issues,
        'checks_run': ['SRM', 'consistency', 'timing']
    }
\end{lstlisting}

\section{Integration Patterns}

\subsection{Feature Flags Integration}

Combine experimentation with feature flag systems.

\lstset{language=Python, caption=Feature Flag with A/B Testing}
\begin{lstlisting}
class FeatureFlagService:
    """Combined feature flags and A/B testing."""

    def __init__(self, experiment_service):
        self.experiment_service = experiment_service
        self.permanent_flags = {}

    def is_enabled(self, flag_name, user_id, default=False):
        """
        Check if feature is enabled for user.

        Supports:
        - Boolean flags (on/off for all users)
        - Percentage rollouts
        - A/B test variants
        """
        # Check permanent flags first
        if flag_name in self.permanent_flags:
            return self.permanent_flags[flag_name]

        # Check if it's an A/B test
        if self.experiment_service.is_experiment(flag_name):
            variant = self.experiment_service.get_variant(flag_name, user_id)
            return variant == 1  # Treat variant 1 as "enabled"

        return default

    def get_variant(self, experiment_name, user_id):
        """Get experiment variant for user."""
        return self.experiment_service.get_variant(experiment_name, user_id)
\end{lstlisting}

\subsection{Metrics Collection}

Integrate with existing analytics systems.

\begin{enumerate}
    \item \textbf{Passive tracking:} Automatic event capture

    \item \textbf{Active instrumentation:} Explicit metric logging

    \item \textbf{Computed metrics:} Derived from raw events

    \item \textbf{Real-time aggregation:} Pre-compute for dashboards
\end{enumerate}

\section{Common Implementation Pitfalls}

\subsection{Assignment Bugs}

\begin{itemize}
    \item \textbf{Issue:} Inconsistent assignments across sessions
    \item \textbf{Cause:} Not using stable user IDs
    \item \textbf{Fix:} Hash-based deterministic assignment
\end{itemize}

\subsection{Survivorship Bias}

\begin{itemize}
    \item \textbf{Issue:} Only analyzing users who complete flows
    \item \textbf{Cause:} Filtering out non-converters
    \item \textbf{Fix:} Intent-to-treat analysis
\end{itemize}

\subsection{Logging Failures}

\begin{itemize}
    \item \textbf{Issue:} Missing assignment or event data
    \item \textbf{Cause:} Client-side logging dropped
    \item \textbf{Fix:} Server-side logging, retry logic, monitoring
\end{itemize}

\section{Summary}

Robust implementation is critical for trustworthy A/B testing:

\begin{itemize}
    \item Use deterministic hash-based randomization
    \item Maintain centralized configuration management
    \item Implement comprehensive logging with quality checks
    \item Integrate with feature flags and analytics
    \item Monitor for common implementation issues
\end{itemize}

Proper infrastructure enables scaling experimentation across the organization.

\section{Further Reading}

\begin{itemize}
    \item \cite{kohavi2020trustworthy}: Implementation best practices
    \item \cite{tang2010overlapping}: Experiment system architecture
    \item Open source platforms: Optimizely, GrowthBook, Unleash
\end{itemize}

\section{Exercises}

\begin{enumerate}
    \item Implement complete hash-based randomization with traffic allocation and salting
    \item Build logging system with SRM detection and automated alerts
    \item Design experiment configuration schema supporting complex targeting
    \item Create feature flag service integrating A/B tests and gradual rollouts
\end{enumerate}
